
We will now evaluate the performance of the IJ covariance on a number of
real-world models and datasets taken from \citet{gelman:2006:arm}.  We limited
our analysis to models for which data was available in the Stan examples
repository \citep{stan-examples:2017}, and for which \texttt{rstanarm}
\citep{rstanarm} was able to run in a reasonable amount of time. Since these are
real-world datasets, the ground truth sampling distribution is not available
from simulations, so, as an imperfect proxy, we compare the IJ covariance to the
(computationally expensive) bootstrap variance instead.

% What are our parameters of interest?
All models we considered from \citet{gelman:2006:arm} were instances of
generalized linear models, including both ordinary and logistic regression as
well as models with only fixed effects and models that include random effects.
Model details, including the number of estimated quantites of each type, are
given in \tabref{arm_models} of \appref{experiment_arm_appendix}.  For
parameters of interest, $g(\theta)$, we took all fixed-effect regression
parameters, the log of the residual variance (when applicable), and the log of
the random effect variance (when applicable). We did not examine random effects,
nor the random effect covariances for the few models with multiple random
effects, out of concern that a BCLT might not be expected to hold for these
parameters.  Below, when we refer to a ``parameter,'' we refer to a scalar
parameter of interest in a particular model for a particular dataset; there were
a total of $\armNumCovsEstimated$ such parameters.

For each model and dataset, we fix a number of MCMC samples, $M =
\armNumMCMCSamples$ to draw for each posterior sample, and a number of
bootstraps samples to draw, $B = \armNumBootstraps$. The total computation time
summed across all models for the initial MCMC and IJ computation was
$\armTotalMCMCTimeMins$ minutes, and the total bootstrap compute time summed
across all models and bootstrap samples was $\armTotalBootTimeHours$
hours.\footnote{Naively, one would expect the total bootstrap compute time to be
roughly $B = \armNumBootstraps$ times the compute time of the single MCMC run,
but the total bootstrap time was in fact longer than this. We conjecture that
the bootstrap time scaled super-linearly due to the fact that some bootstrap
samples led to posteriors that were less well-conditioned than the original
datasets, and so incurred more rejections during the MCMC procedure.}
%
For each parameter of interest, we use a single run of MCMC
to compute the IJ covariance estimate, $\gcovijhat$.  Using \algrref{boot}, we
draw $B$ different bootstrap samples, compute $M$ MCMC samples for each, and
finally compute the bootstrap variance estimate, $\gcovboothat$.

All together, this procedure produces comparisons between $\armNumCovsEstimated$
distinct covariance estimates from $\armNumModels$ different models using
$\armNumDatasets$ distinct datasets. The median number of observations amongst
the models was $\armMedianNumObs$, and the range was from $\armMinNumObs$ to
$\armMaxNumObs$ exchangeable observations.  Since we fit our models in
\texttt{rstanarm}, we were able to use the function \texttt{rstanarm::log\_lik}
to compute $\ell(\x_n \vert \theta)$, so $\gcovijhat$ required essentially no
additional coding.  We used the default priors provided by \texttt{rstanarm}.




\ARMGraphDiff{}


\ARMGraphZ{}


% Why do we need summary metrics
With $\armNumCovsEstimated$ covariances to report, we must summarize our results
into a few model--agnostic high-level accuracy summaries.  To this end, for each
parameter $j$ in each model $i$, we computed the following measure of both the
``practical significance'' and ``statistical significance'' of the IJ covariance
covariance as a proxy for the bootstrap variance:
%
\begin{align}
%
\normdiffij_{ij} :=
     \frac{\gcovijhat_{ij} - \gcovboothat_{ij}}
     {\left| \gcovboothat_{ij} \right| + \seboot_{ij}}
\quad\textrm{and}\quad
\zdiff_{ij} :=
\frac{\gcovijhat_{ij} - \gcovboothat_{ij}}
     {\sqrt{(\seij_{ij})^2 + (\seboot_{ij})^2}}.
\eqlabel{arm_metrics}
%
\end{align}
%
In \eqref{arm_metrics}, the quantities $\seij$ and $\seboot$ are measures of
the Monte Carlo error in the corresponding covariance estimates
(see \appref{experiments_sampling_error} for details).

The metric $\normdiffij_{ij}$ is a (noisy) measure of the proportional
difference between $\gcovijhat_{ij}$ and $\gcovboothat_{ij}$. The term
$\seboot_{ij}$ is included to retain a meaningful notion of scale when
$\gcovboothat_{ij}$ is nearly zero due to chance alone.
%
The metric $\zdiff_{ij}$ is the z-statistic which we would compute to test
equality between $\gcovijhat_{ij}$ and $\gcovboot_{ij}$ using our estimates of
the sampling error.  Despite this interpretation, the z-statistics for different
parameters within a given model are based on the same MCMC samples and so are
not independent, and we should be cautious in interpreting the collection of
$\zdiff_{ij}$ as the basis of any kind of formal hypothesis test. Rather, we
recommend thinking of $\zdiff_{ij}$ as heuristics whose marginal distributions
summarize the error $\gcovijhat - \gcovboot$ relative to Monte Carlo error.


% \footnote{Note that the relative errors $\normdiffij$ and
% $\normdiffbayes$ do not attempt to remove or account for Monte Carlo error.  One
% might imagine computing the relative error after deconvolving the Monte Carlo
% error and covariance estimates using, say, nonparametric empirical Bayes.  We
% investigated such an approach, but the resulting procedure was complicated and
% the results qualitatively similar to \figref{normerr_graph} below, so we will
% report only the relatively simple metrics $\normdiffij$ and $\normdiffbayes$.}
%
The distributions of $\normdiffij_{ij}$ and $\zdiff_{ij}$ across models,
datasets, and parameters are shown in \figref{normerr_graph,relerr_graph}
respectively.
%
For models with more data, the IJ provides a reasonable approximation to the
much more computationally expensive bootstrap, with roughly half of the IJ
covariance estimates falling well within 50\% of the bootstrap covariance
estimates, and rejecting the hypothesis test of equivalence at roughly the
nominal rate.  For models with less data, IJ estimates do not appear to
systematically deviate from the bootstrap estimates, but are more variable.
Recall that we expect, theoretically, only that the IJ and bootstrap variances
coincide as $N \rightarrow \infty$.

The distribution of $\zdiff_{ij}$ suggests that the apparent differences between
the IJ and bootstrap can be accounted for by Monte Carlo error alone in the
large models but not necessarily the small ones.  Recalling that both the
bootstrap and the IJ covariances are distinct approximations of the same
quantity which are close to one another only asymptotically, we conjecture that
the needed asymptotics have not yet kicked in for models with fewer datapoints.

% In \appref{experiment_arm_appendix}, we show that the Bayesian posterior
% covariance is a poor approximation for the bootstrap in larger datasets,
% particularly for covariances involving at least one scale parameter. We
% conjecture that the gap between the IJ covariance and the posterior covariance
% for log scale parameters may be driven by misspecification, though a more
% careful investigation is beyond the scope of the present work.

